%
% Hauptdokument
% Abschlussarbeit (Bachelor)
%
% Thema: Erstellung einer Browser Extension zur Usability Evaluierung von beliebigen Web-Applikationen �ber Heatmaps.
% Betreuer 1: Prof. Dr. Targo Pavlista
% Betreuer 2: Siamak Haschemi
%
% @author Christian Bromann <contact@christian-bromann.com>
%

\documentclass[
	bibliography=totoc,
	listof=totoc,					% falls Index verwendet, erg�nze "index=totoc" zu den Optionen 
	BCOR=5mm,					% Rand f�r Bindung: 5mm
	DIV=12]{scrbook}
\usepackage{bibgerm}       			% deutsche Literaturverzeichnisse
\usepackage[applemac]{inputenc} 		% Umlaute im Text
\usepackage{graphicx} 				% Einf�gen von Grafiken  - f�r PDF-Latex: .pdf und .png (.jpg m�glich, sollte aber vermieden werden)

\usepackage{fancyvrb, listings, color} 	% F�r Code Ausgaben mit Highlighting
\usepackage{caption}

\definecolor{grey}{RGB}{240,240,240}
\definecolor{darkgrey}{RGB}{88,88,88}
\definecolor{brown}{RGB}{237,201,175}
\definecolor{darkbrown}{RGB}{146,111,91}
\definecolor{darkgreen}{RGB}{0,100,50}
\definecolor{darkpurple}{RGB}{75,0,130}
\definecolor{php}{RGB}{192, 64, 0}
\definecolor{vars}{RGB}{71, 110, 155}
\definecolor{functions}{RGB}{156, 112, 63}

\lstset{
    fancyvrb=true,
    showstringspaces=false,
    tabsize=3,
    keywordstyle=\color{blue}\textbf,
    commentstyle=\color{darkgreen}\textit,
    stringstyle=\color{red},
    backgroundcolor=\color{grey},
    emph=[1]{php},
    emphstyle=[1]\color{php},
    emph=[2]{public,function,use,as,class,private,namespace,return},
    emphstyle=[2]\color{darkpurple},
    emph=[3]{Signee},
    emphstyle=[3]\color{darkbrown},
    emph=[4]{null,array},
    emphstyle=[4]\color{brown},
    emph=[5]{$form,$request,$em,$this,$id,$createdAt,$name,$signee,$signees},
    emphstyle=[5]\color{vars},
    emph=[6]{getSigneeForm,getRequest,getMethod,getData,getDoctrine,getEntityManager,persist,flush,getRepository,render,findAll},
    emphstyle=[6]\color{functions}
}

\DeclareCaptionFont{white}{\color{white}}
\DeclareCaptionFormat{listing}{\colorbox{darkgrey}{\parbox{\textwidth}{#1#2#3}}}
\captionsetup[lstlisting]{format=listing,labelfont=white,textfont=white}

\usepackage{url}           				% URL's (z.B. in Literatur) sch�ner formatieren
\usepackage{hyperref} 				% sorgt f�r Hyperlinks in PDF-Dokumenten

\graphicspath{{./images/}}				% Pfad zu den Bildern des Dokumentes

\begin{document}

% ---------------------------------------------------------------

\frontmatter 						% Titelbl�tter und Erkl�rung
    %
% Hauptdokument
% Praktikumsbericht zur Anerkennung der Arbeit bei der Spot-Media AG
% @author Christian Bromann <contact@christian-bromann.com>
%

\thispagestyle{empty}

\begin{center}

\vspace*{-1cm}

\includegraphics[width=0.3\textwidth]{beuth_logo}

\vspace{2cm}

{\Large \textbf{Bachelorarbeit}}\\ 

\vspace{1cm}

{\Huge \textbf{Erstellung einer Browser Extension}}\\
\vspace*{3mm}
{\Huge \textbf{zur Usability Evaluierung von}}\\
\vspace*{3mm}
{\Huge \textbf{beliebigen Web-Applikationen �ber}}\\
\vspace*{3mm}
{\Huge \textbf{Heatmaps}}\\

\vspace{2cm}

{\Large \textbf{Christian Bromann}}\\ 
\textbf{Matr. Nummer 774261}

\vspace{2cm}

\parbox{1cm}{
\begin{large}
\begin{tabbing}
Betreuer: \hspace{1.5cm}\=Prof. Dr. Targo Pavlista\\[1mm]
\>Siamak Haschemi\\[2mm]
Gutachter: \>Max Mustermann\\[2mm]
Abgabetermin: \> 7. Juli 2013\\
\end{tabbing}
\end{large}}\\

Draft vom \today 			% TODO: Entfernen bei Abgabe

\end{center}

\newpage
 				% Titelblatt mit Beuth Logo, Titel und Name
    \thispagestyle{empty}
    \cleardoublepage
    %
% Erklärung zur selbstständigen Arbeit
% Praktikumsbericht zur Anerkennung der Arbeit bei der Spot-Media AG
% @author Christian Bromann <contact@christian-bromann.com>
%

\newpage

\thispagestyle{empty}

\begin{large}

\vspace*{1cm}

\noindent
{\Huge \textbf{Ehrenwörtliche Erklärung}}\\
\\
Hiermit erkläre ich, Christian Bromann, geboren am 04.08.1989 in Salzwedel, ehrenwörtlich, dass ich meine Bachelorarbeit mit dem Titel:\\
\\
\textit{„Erstellung einer Browser Extension zur Usability Evaluierung von beliebigen Webapplikationen über Heatmaps“}\\
\\
selbstständig und ohne fremde Hilfe angefertigt und keine anderen als in der Abhandlung angegebenen Hilfen benutzt habe.
Die Übernahme wörtlicher Zitate aus der Literatur sowie die Verwendung der Gedanken anderer Autoren an den entsprechenden Stellen habe ich innerhalb der Arbeit gekennzeichnet.\\
\\
Ich bin mir bewusst, dass eine falsche Erklärung rechtliche Folgen haben kann.

\vspace{2cm}

\noindent
Berlin, den \today

\vspace{3cm}

\hspace*{7cm}%
\dotfill\\
\hspace*{9.5cm}%
\textit{Christian Bromann}

\end{large}
		% Erkl�rung zur selbstst�ndigen Arbeit
    \vspace*{2cm}

\begin{center}
    \textbf{Abstract}
\end{center}

\vspace*{1cm}

\noindent Hier steht eine kurze Zusammenfassung der Arbeit. Sie darf auf gar keinen Fall
länger als eine Seite sein, ca. eine drittel bis eine halbe Seite ist optimal.

 			% Abstract
    \thispagestyle{empty}
    \tableofcontents 					% Inhaltsverzeichnis

% ---------------------------------------------------------------

\mainmatter 						% die eigentliche Arbeit

    \chapter{Einleitung}
	\section{Motivation und Ausgangssituation}
	\section{Projektziel}
	\section{Aufbau der Arbeit}

\chapter{Web-Usability}
	\section{Normen}
	\section{Barrierefreiheit}
	
\chapter{thEvaluator}
	\section{Zielbestimmungen}
	\section{Produkteinsatz}
	\section{Produktfunktion}
	\section{TestCase Definition}
	\section{Aufzeichnungskriterien}
	\section{Qualit�tsanforderungen}
	\section{Entwicklungsumgebung}

\chapter{Datenverarbeitung}
	\section{Sockets}
	\section{NoSQL Datenbank}
	\section{Datenaufbereitung}

\chapter{Pr�sentation}
	\section{Heatmaps}
	\section{Kriterien}
	\section{WebGL Ausgabe}




% ---------------------------------------------------------------

\backmatter 						% ab hier keine Nummerierung mehr
    \listoffigures						% Abbildungsverzeichnis
    \lstlistoflistings					% Listing
    %\include{./Literaturverzeichnis/literaturverzeichnis}

\end{document}
